\documentclass[a4paper,10pt]{scrartcl}
%\documentclass[a4paper,10pt]{scrartcl}

\usepackage[xetex]{graphicx}
\usepackage{fontspec}

\title{CICO}
\author{HAUM}
\date{26/02/2016}

\begin{document}
\maketitle

\hfill
\raisebox{-.5\height}{\includegraphics[width=6cm]{logo}}

\begin{minipage}{10cm}
Organisme/Structure : HAUM\\
Adresse : 19 boulevard Marie \& Alexandre Oyon\\
Prénom, Nom contact : Romuald Conty\\
Titre : Président\\
Tel portable : 06.60.09.86.20\\
Mail : romuald@haum.org\\
\end{minipage}


\section{Présentation de vos activités et en particulier par rapport à l’innovation}

Le HAUM est une association sarthoise qui structure un Hackerspace AU Mans, littéralement un lieu pour bidouilleurs.

En tant que tel, l'association anime un lieu de vie, un endroit où des passionnés, des curieux, des experts, etc. se retrouvent, collaborent et échangent. C'est également un lieu de découverte, d'apprentissage, d'expérimentation, de discussion et de partage.

En favorisant ainsi les rencontres de personnes, quelque soit leur bagage technique ou leurs objectifs, et en soutenant la transmission de savoirs et de savoir-faire, l'association a su attirer des profils variés dont les échanges ont mené à des projets innovants.

Aujourd'hui le HAUM dispose d'une expérience réelle pour aider à capter et cristalliser des idées innovantes dont les réalisations effectives sont régulièrement plébiscitées. La présence d'un lieu facilement accessible, disposant de matériel partagé, et favorisant les échanges entre publics hétéroclites (entrepreneurs, artistes, techniciens, passionnés...) semble important pour stimuler la créativité et l'innovation sur Le Mans. 

\subsection*{Fabrication attractive}
À titre d'exemple en 2014, le HAUM réalise un jeu lumineux (Pong1D) pour l'évènement culturel "les Siestes Teriaki" à l'abbaye de l'Épau. Le projet séduit le public par son caractère atypique et la qualité de la réalisation. L'association est désormais très sollicitée pour présenter ce projet, ce qui a permis par exemple de nouer des liens avec le fablab de Nantes (porté par l'association Ping) qui a demandé une participation du HAUM à leur "FestivalD", ainsi qu'avec d'autres hackerspaces et fablabs.

\subsection*{Expertises libérées}
Aussi, la diversité des sensibilités et compétences présentes au sein de l'association, ainsi que les libertés totales accordées aux acteurs de la création, ont permis de concrétiser des projets qui n'auraient pas pu exister dans des structures plus codifiées.
Pour exemple, le dHaum -- un instrument de musique géant dont les visiteurs, en se déplaçant dans l’installation, constituent à la fois les musiciens et des éléments de l'instrument lui-même -- a émergé grâce à cette souplesse.

\subsection*{Actions collaboratives}
Dans son action, l'association s'investit également régulièrement avec d'autres structures partageant l'envie de "faire" comme Teriaki pour des animations au cours des Siestes 2014 puis du Festival 2015, Ping pour le FestivalD à Nantes, la Ruche Numérique avec des actions de formations techniques, de présentation d'outils, de conférences, etc., l'ENSIM avec l'élaboration et l'animation de sujets pour les 24h du code, les conseils de quartiers avec une coorganisation de Repair Cafés, les jeudis du libre avec Linux Maine, la journée de conférences AgileMans, et bien d'autres.

\subsection*{Aventure intellectuelle}
En outre, le hackerspace ne se limite pas à des activités technologiques et explore des domaines contrastés venant enrichir le champ d'innovation. Indiquons par exemple le recyclage de cartons usagés pour la réalisation de mobilier, l'exploration de topographies mathématiques ayant donné lieu à la création de casses-têtes ingénieux, diverses expérimentations dans les domaines du design et de la photographie, etc.

\subsection*{Savoirs partagés}
Le HAUM fournit un environnement accueillant pour apprendre, découvrir et partager.
Dans ce contexte, pour favoriser la diffusion des connaissances, le hackerspace organise des cycles de conférences courtes pour que chacun puisse présenter un sujet qui lui tient à cœur et que tous puissent découvrir des sujets en dehors de leurs intérêts directs. De même, pour les sujets demandant plus de temps d'appréhension, l'association a mis en place des sessions de formation et à l'occasion des ateliers.

\section{Vos idées et commentaire sur l’utilisation des deux espaces repérés sur Novaxud}
\emph{(Bâtiment ex « chaufferie » de 380 m² et surface bureaux disponibles de 1 000 m²)}

Le HAUM est très favorable à l'utilisation de l'ex chaufferie.

Il est important de ne pas figer son aménagement et au contraire livrer un lieu brut, modulable, qui pourra être reconfiguré au gré des besoins dictés par les innovations. Tantôt lieu de fabrication, tantôt lieu d'exposition, il faut que le lieu puisse vivre et s'adapter aux utilisateurs de l'espace, pour des projets ponctuels ou de plus longue durée. 

Plus qu'un simple espace de travail, le fablab doit être un lieu de vie où chacun se sent chez lui, dans un environnement propice à l'émergence de nouvelles idées.
Il renferme tous les ingrédients nécessaires pour passer d'une idée dans la tête d'un porteur de projet à un prototype sur le bureau du financeur.

Les membres du HAUM sont prêts à aider au démantèlement et souhaitent participer à son aménagement avec l'ensemble des volontaires, afin de s'assurer de créer un lieu de vie à l'image de ses occupants.

D'un point de vue pratique, il serait utile de prévoir une grande entrée afin d'y introduire des machines volumineuses et souvent nécessaires dans de tels lieux. En effet, certains fablabs ont déjà été contraints d'effectuer de la maçonnerie pour installer des machines. Le prévoir dès le début du projet peut éviter les déconvenues.

En ce qui concerne l'espace de 1000m², il a l'avantage d'être déjà disponible pour débuter le projet de cette \emph{Cité de l'Innovation Collaborative}. Rapidement, il serait cependant préférable de le remplacer par les deux bâtiments à coté de la chaufferie, plus propices au développement du projet, et offrant un cachet, une atmosphère propre à la Cité. 

% \section{Résumé l'activité du HAUM}
% 
% Le HAUM est une association sarthoise créée il y a 3 ans qui axe son action autour de la transmission de savoirs et de savoir-faire. Cette structure propose à chacun, quelque soit son bagage technique, de venir rencontrer d'autres passionnés, par exemple lors de séances ouvertes au public
% 
% Depuis 2 ans, l'activité de l'association a pris de l'ampleur. Elle s'adresse à une audience assez large au travers de projets accessibles aux néophytes : pong 1D (jeu), PianoStairs et dHAUM (instruments de musique géants et collaboratifs), meubles en carton, etc.
% 
% Dans son action, le HAUM s'associe régulièrement avec des structures partageant l'envie de "faire" : Teriaki (pour les siestes 2014 puis le festival 2015), Ping (pour FestivalD, à Nantes), la Ruche Numérique (hébergement, formations, conférences, etc.), ENSIM (24h du code), Conseils de Quartiers (Repair Cafés) et continuera : Gamer Assembly 2016 (Poitiers), etc.
% 
% Forte de ces interactions, pleine de projets variés et originaux, l'association compte aujourd'hui une trentaine de membres et poursuit toujours ses deux objectifs principaux. Le premier est de fournir à tous un environnement accueillant pour apprendre, découvrir et partager. Cet objectif s'articule autour de projets communs qui avancent avec la communauté ou de projets personnels ayant besoin d'aide extérieure.  
% 
% Le second objectif est la diffusion des connaissances. À ce titre, le HAUM organise des cycles de conférences courtes pour que chacun puisse présenter un sujet qui lui tient à cœur et que tous puissent découvrir des sujets en dehors de leur intérêt direct. De même, pour les sujets demandant plus de temps d'appréhension, l'association a mis en place des sessions de formation et à l'occasion des ateliers.
% 
% Tous ces points rejoignent un idéal d'auto-formation, d'apprentissage des techniques au travers de l'expérimentation et au travers de la rencontre avec d'autres. Cet aspect est essentiel dans un projet global de \emph{\emph{Cité de l'Innovation Collaborative} Collaborative} car il apporte le renouveau permanent et la richesse d'idée nécessaire à tout projet en gestation.
% 

% \section{Qu'est ce qu'un fablab ? (selon le HAUM)}
% 
% Un fablab, au delà des seuls intérêts matériels qu'il représente (accès à des machines, des outils, etc.), est le lieu commun qui, modifiable par essence, ressemble aux personnes qui l'utilisent.
% 
% Un projet développé dans un fablab sera continuellement soumis à la critique : c'est un excellent moyen d'avancer rapidement. Le fablab doit aussi être un endroit qui fourmille d'idées s'il veut développer, vivre et faire vivre l'innovation au sein de la Cité. Cela n'est évidemment pas incompatible avec l'exploitation commerciale pour le porteur de projet. 
% 
% Dans le cadre des actions du fablab, il y a bien évidement l'organisation d’événements à destination du public : conférences, ateliers, démonstrations, etc.
% 
\section{Que manque-t-il au HAUM pour offrir le meilleur à une cité de l'innovation ?}

\subsection*{Extension des horaires d'ouvertures}

L'intérêt d'une association comme le HAUM dans le projet de la \emph{Cité de l'Innovation Collaborative}, et en particulier sur l'axe fablab, est la communauté qui gravite autour.
Afin d'offrir un accompagnement efficace aux créateurs, le HAUM a besoin d'étendre ses plages d'ouverture. Aujourd'hui, les activités de chacun nous permettent d'ouvrir seulement le mardi et le jeudi soir de 18h30 à 1h (parfois plus). Ponctuellement, nous ouvrons le lieu le samedi, dimanche ou mercredi après-midi.
Cependant l'accueil bénévole n'est pas suffisant pour tous les publics et doit être complété pour correspondre aux besoins de l'ensemble des acteurs ; idéalement un lieu ouvert de manière quasi-permanente.

\subsection*{Matériel}

Pour profiter pleinement du lieu, les acteurs auront besoin d'avoir accès à des équipements de création.
Sans être exhaustif, un fablab s'équipe généralement du matériel suivant :

\begin{itemize}
 \item Imprimante 3D
 \item Usinage à commande numérique (fraiseuse, tour)
 \item Découpeuse laser
 \item Découpeuse vinyle
 \item Travail textile à commande numérique (tricoteuse, brodeuse)
 \item Équipement de base pour l'électronique (PCB, station de soudage, mesure)
\end{itemize}

Il s'agit de mutualiser les moyens pour prototyper, fabriquer, tester, concrétiser, désassembler, modifier... et ainsi faire en sorte que tout créateur puisse accéder aux moyens habituellement réservés à l'industrie.

\textit{Note :} au besoin, le HAUM est volontaire pour contribuer à l'élaboration d'une liste plus complète.

\subsection*{Lieu}

Au delà de l'aspect fabrication, différents espaces sont necessaires pour permettre une meilleure collaboration :
\begin{itemize}
 \item salle de conférence
 \item salle de réunion
 \item paillasses / établis
 \item espace détente / restauration
\end{itemize}

\textit{Note :} au besoin, le HAUM est volontaire pour approfondir les aspects pratiques et apporter son expérience à l'aménagement d'un tel lieu.

\section{Comment anime-t-on un fablab ?}
L'animation d'un fablab est un point crucial, en particulier sur les premiers mois/années. En effet, toute la puissance d'un espace de rencontre, comme celui de la \emph{Cité de l'Innovation Collaborative}, réside dans la communauté qui l'habite, l'anime et le transforme.
Le fablab s'articule autour d'une communauté qui s'améliore au cours du temps, il faut donc un élément premier, fédérateur de la communauté naissante.

Idéalement il y a deux axes majeurs dans l'animation d'un fablab : l'axe technique et l'axe communication.

La technique, tout d'abord, est l'objet, la raison d'être du fablab. Le fablab étant un lieu où chacun vient pour réaliser des objets, il est bien sûr nécessaire que l'équipe animatrice soit compétente.
La gestion d'un fablab comporte des points amusants, passionnants (projets, rencontres, formations, etc.) mais aussi des éléments plus complexes : gestion du stock de matières premières, maintenance/entretien des machines, relations extérieures, gestion des pannes/de l'utilisation des machines, etc.
Si elle ne requiert pas forcément une formation spécialisée, cette deuxième facette du \textit{fab manager} (personne gérant le fablab) demande toutefois un investissement qu'on ne saurait attendre d'un bénévole.

Les animations à mettre en place du côté technique sont nombreuses : ateliers et formations sur les machines, mini-projets en groupe pour se former à certaines technologies, conférences sur divers sujets (électronique, matériaux, techniques de prototypage, etc.) mais aussi défis, \textit{challenges} et autres activités permettant de transformer un ensemble d'usagers en une communauté soudée.

Il ne faut jamais perdre de vue que, au delà de sa mission technique, le \textit{fab manager} cristallise l'âme du lieu : il re-motive ceux dont le projet patine, il aide de manière désintéressé quiconque le demande, il gère les éventuels conflits de personne et sert de médiateur en cas de besoin. L'aspect social du \textit{fab manager} est essentiel et parmis tous les critères possibles, c'est celui qu'il faut mettre en avant dans un recrutement : on peut apprendre à se servir de machines mais on apprend que difficilement à composer avec l'aspect social d'une communauté partageant un lieu à longueur de journée.

Le deuxième axe de l'animation d'un fablab concerne la communication : le fablab prend de la valeur au travers de sa communauté et des compétences qui se développent.

Une bonne communication, c'est l'assurance d'une entrée continue de nouvelles personnes, c'est l'assurance d'entrées dont on connait a priori une partie des envies. Une bonne communication c'est aussi une grande visibilité pour les occupants : article(s) sur leur projet, médiatisation, mise en avant et \textit{showcase}, autant de choses essentielles pour la reconnaissance dont chaque porteur de projet a besoin.

Pour résumer, une animation de fablab requiert une ou deux personnes (la technique étant un poste à temps complet) qui se répartissent des tâches de communication et d'aide aux occupants. Le fablab doit s'appuyer pour grandir sur un ensemble cohérent d’événements plus ou moins médiatiques axés sur les formations/conférences proposées mais aussi et surtout sur les projets en cours au fablab. Dans les activités de communication, il est essentiel de participer à des événements (en Sarthe ou ailleurs) organisés par d'autres structures : dans la culture fablab, la reconnaissance vient des autres communautés.

\end{document}
